\section{实验原理}
\subsection{SM4算法}
SM4算法采用非平衡Feistel结构,分组长度为128bit,密钥长度为128bit,迭代轮数为32轮。
\subsubsection{加密算法}
输入128bit明文\(m\),按32-bit分为4个字\(X^0_0,X^1_0,X^2_0,X^3_0\)
进行32轮迭代,每轮迭代都有一个32bit的轮密钥$K_{i-1}$参与计算,第$i(i = 1,2,\cdots,32)$轮的计算如下:
\[
\begin{aligned}
    X^0_{i} &= X^1_{i-1} \\
    X^1_{i} &= X^2_{i-1} \\
    X^2_{i} &= X^3_{i-1} \\
    X^3_{i} &= X^0_{i-1} \oplus T(X^1_{i-1}\oplus X^2_{i-1}\oplus X^3_{i-1}\oplus K_{i-1})
\end{aligned}
\]
对最后一轮迭代的输出进行反序变化$R$得到128bit密文\(c\):
\[
c = (Y^0,Y^1,Y^2,Y^3) = R(X^{0}_{32},X_{32}^1,X_{32}^2,X_{32}^3) = (X^3_{32},X^2_{32},X^1_{32},X^0_{32})
\]
其中函数$T:\mathbb{F}_2^{32}\to \mathbb{F}_2^{32}$由非线性变换\(S\)和线性变换\(L\)组成:
\[    T = L\circ  S   \]
设函数$T$的32bit输入为$A = (a_0,a_1,a_2,a_3) \in (\mathbb{F}_{2^8})^4$,则
\begin{itemize}
    \item 非线性变换$S$:对每个字节单独进行查表代替操作,即
    \[ B = (S(a_0),S(a_1),S(a_2),S(a_3))  \]
类似AES加密,SM4算法采用1个S盒,查表时将每个字节的高4bit看作行标,低4bit看作列标,对应取值即为输出字节。
    \item 线性变换$L$:对非线性变换$S$的32bit输出$B$,计算
    \[   L(B) = B \oplus (B \lll 2) \oplus (B \lll 10) \oplus (B \lll 18) \oplus (B \lll 24) \]
    其中$\lll x$表示循环左移$x$bit。
\end{itemize}
\subsubsection{轮密钥生成方案}
SM4算法的主密钥$K$为128bit,轮密钥生成方案也采用非平衡Feistel结构。
轮密钥\(K_i\)由主密钥\(K\)通过密钥扩展算法生成,具体做法为:
先将主密钥\(K\)按32-bit分为4个字\(MK_0,MK_1,MK_2,MK_3\),然后与32-bit常量\(FK_i\)进行异或得到初始轮密钥:
\[
\begin{aligned}
    K_{-4} &= MK_0 \oplus FK_0 \\
    K_{-3} &= MK_1 \oplus FK_1 \\
    K_{-2} &= MK_2 \oplus FK_2 \\
    K_{-1} &= MK_3 \oplus FK_3
\end{aligned}
\]
其中$FK_0 = $0xA3B1BAC6,$FK_1 = $0x56AA3350,$FK_2 = $0x677D9197,$FK_3 = $0xB27022DC。
接着通过迭代计算得到后续轮密钥:
\[
K_i = K_{i-4} \oplus T'(K_{i-3} \oplus K_{i-2} \oplus K_{i-1} \oplus CK_i), \quad i=0,1,...,31
\]
其中$CK_i$为固定参数;\(T'\)函数与\(T\)函数类似,但线性变换部分不同,仅需替换$L$为$L'$:
\[
\begin{aligned}
    T'(X) &= L'(S(X)) \\
    L'(B) &= B \oplus (B <<< 13) \oplus (B <<< 23)
\end{aligned}  
\]

\subsubsection{解密算法}
因为采用广义Feistel结构,所以SM4算法具有加解密结构一致性,因此解密算法与加密算法一致,仅需将参与第$i(i = 1,2,\cdots,32)$轮迭代
运算的轮密钥由$K_{i-1}$替换为$K_{32-i}$即可。




\subsection{CBC模式}
SM4只能加密128bit明文信息,为了能实现加密任意长度的明文信息$m$,需要先将明文填充至长度为128的整数倍,再将填充后的结果换分为多个块大小为128bit的待加密明文块
    \[ P = P_1 || P_2 || \cdots P_s,P_i \in\{0,1\}^{128} \]
CBC模式是一种分组加密模式,它将每个明文块与前一个密文块进行异或操作后再进行加密,以增强安全性。
\subsubsection{CBC加密}
CBC模式的加密过程如下:
\[
\begin{aligned}
    C_0 &= E(P_0 \oplus IV) \\
    C_i &= E(P_i \oplus C_{i-1}), \quad i \geq 1
\end{aligned}
\]
其中\(P_i\)是第\(i\)个明文块,\(C_i\)是第\(i\)个密文块,\(IV\)是初始化
向量,必须是随机的且唯一的。CBC模式通过引入IV和前一个密文块的依赖关系,使得相同的明文块在不同的位置产生不同的密文块,从而提高了加密的安全性。也正因为当前密文块依赖前一密文块,CBC加密过程在单条消息内部不能直接按分组并行。
\subsubsection{CBC解密}
CBC模式解密过程与加密过程类似,需要用上一个密文分组与本组密文分组解密后的结果异或,得到最终明文,公式如下:
\[
\begin{aligned}
    P_0 &= D(C_0) \oplus IV \\
    P_i &= D(C_i) \oplus C_{i-1}, \quad i \geq 1
\end{aligned}
\]
虽然CBC模式解密时当前明文分组需要与前一分组密文异或,但实际解密时所有密文分组都已知,因此各分组的\(D(C_i)\)可以提前并行计算,再分别与对应的前序密文分组或IV异或得到明文。

\subsection{SIMD指令}
SIMD(Single Introduction Multiple Data)即单指令流多数据流,是一种采用一个控制器来控制多个处理器,同时对一组数据(数据向量)中的每一个分别执行相同
的操作,从而实现空间上的并行性的技术,即一个指令能够同时处理多个数据。



\subsection{RISC-V密码学拓展}
RISC-V密码学拓展(K拓展)主要为解决软件密码学的性能低下和侧信道攻击风险的问题,分为标量密码学和向量密码学两大体系。
\begin{itemize}
    \item [1.]标量密码学拓展:主要面向轻量级嵌入式设备或不包含向量单元的通用处理器。它们复用CPU原有的32个通用标量寄存器。
    \begin{itemize}
        \item Zkne:AES加密指令,支持32位和64位寄存器操作。
        \item Zksh:SM3哈希指令,硬件实现压缩函数。
        \item Zkr:熵源拓展,提供硬件随机数生成。
        \item Zksh:SM4加速指令。
    \end{itemize}
    \item [2.]向量密码学拓展:主要面向高性能服务器、数据中心等高吞吐场景。它将密码学算法与RISC-V向量拓展结合。
\end{itemize}
RISC-V B拓展部分实现了特殊位操作和变换,使用专门的指令实现这些操作能大幅提升效率。
\begin{itemize}
    \item 循环移位:ror,rol。
    \item 位或字节排列组合:rev.b,rev8
    \item 无进位乘法:clmul,clmulh
    \item 带取反逻辑:andn,orn,xnor
    \item 装箱:pack,packu,packh
    \item Crossbar组合:xperm.n,xperm.b
\end{itemize}


\subsection{SM4 bit slice}
\footnote{参考GmSSL:\url{http://gmssl.org/news/20180528.html}}
由于SM4算法在实现过程中常需要根据内部状态执行查表操作,查表的访问模式可能泄露与输入相关的信息,从而使基于Cache时间的侧信道攻击成为可能。SM4 bit slice的基本思想是将查表替换为布尔逻辑门运算,并将多条独立数据流的对应比特位打包到同一组寄存器中并行处理,以获得更接近常量时间的执行特性。
\subsubsection{原理说明}
\begin{itemize}
    \item 算法结构：算法将一个字节的信息按照高位4bit，低位4bit分开存储，因此1个64寄存器可以保存16个高位或者16个低位信息，而一个32位字可以分解成4个高位和4个低位，分开保存在8个64位寄存器内。算法的整体结构基于这种存储结构对标准的SM4算法进行改造，保证单个SM4加密在算法全过程中只使用4bit宽度的存储位，从而在64位寄存器中可以实现16路SM4并行。
    \item $SBOX$拆分实现：SM4的S盒由矩阵运算和8次多项式构成的有限域求逆元运算构成，S盒的一般实现是根据算法的内部状态值，查找内存中表的对应值来实现的，不同的状态值查找的时间可能不同，这也是cache攻击能够产生效果的原因。为了使S盒的查找符合常量时间要求，并满足bit slice算法的结构，对S盒进行拆分，查找过程是输入状态值的4bit高位和4bit低位，输出也是4bit高位和4bit低位。查表算法中，8bit矩阵运算被分解成4bit内的矩阵运算，而8次多项式构成的有限域求逆，则按照复合域同构映射的方法，先映射到有限域
$GF(4^2)$上，再映射到有限域$GF(((2)^2)^2)$上，使得$GF(8)$上的有限域求逆过程，转化成
$GF(((2)^2)^2)$上的运算，其求逆运算被降低到了$GF(2)$上的求逆（可以用简单的函数实现）。经过此方法的分解，S盒的查表运算变成一个可以在寄存器内通过布尔逻辑实现、且运行时间与输入无关的bit slice算法。其实际并行路数取决于具体后端使用的寄存器宽度和平台实现方式。
    \item 并行路数并不是固定常数,而是取决于具体实现所使用的寄存器宽度和后端结构。若仅使用64位通用寄存器,可支持较低路数的并行;若使用更宽的向量寄存器,则可以承载更多独立数据流。因而bit slice更适合作为一种实现方法论,其具体并行度应结合目标平台实际说明。
\end{itemize}








\subsection{性能分析}
对于第 \(i\) 轮，首先需要计算
\[
X_{i+1}\oplus X_{i+2}\oplus X_{i+3}\oplus rk_i
\]
这一部分大约需要 3 次 32 位异或操作。然后进入 S 盒替换。由于输入是 32 位，可以拆分为 4 个字节：
\[
A=(a_0,a_1,a_2,a_3)
\]
每个字节经过一次 S 盒替换，因此每轮包含
\[
4 \text{ 次 S 盒查表}
\]
S 盒替换之后，需要进行线性变换：
\[
L(B)
=
B\oplus (B\lll 2)\oplus (B\lll 10)
\oplus (B\lll 18)\oplus (B\lll 24)
\]
其中，\(\lll\) 表示循环左移。该线性变换大约需要 4 次循环左移和 4 次 32 位异或操作。
最后，还需要与 \(X_i\) 进行一次异或：
\[
X_{i+4}=X_i\oplus L(B)
\]
因此每轮约
\[
8\text{ 次 32位异或}
+
4\text{ 次循环移位}
+
4\text{ 次S盒查表}
\]
因此，单个 128 bit 分组的操作量可以近似表示为
\[
\boxed{
T_{\text{block}}
\approx
256XOR + 128ROT + 128SBOX
}
\]
其中，\(XOR\) 表示 32 位异或操作，\(ROT\) 表示 32 位循环移位操作，\(SBOX\) 表示 8 位 S 盒查表操作。相比单纯地说单分组复杂度为 \(O(1)\)，这种分析方式能够更清楚地体现 SM4 的实际计算开销。

\subsubsection{密钥扩展复杂度}

SM4 密钥扩展过程也需要生成 32 个轮密钥。轮密钥递推公式为

\[
K_{i+4}
=
K_i \oplus T'(K_{i+1}\oplus K_{i+2}\oplus K_{i+3}\oplus CK_i)
\]
其中
\[
T'=L'\circ S
\]
密钥扩展中的线性变换为
\[
L'(B)=B\oplus (B\lll 13)\oplus (B\lll 23)
\]
因此，密钥扩展每轮大致包含：
\[
6 \text{ 次异或}
+
2\text{ 次循环移位}
+
4\text{ 次S盒查表}
\]
32 轮密钥扩展总操作量约为
\[
\boxed{
192XOR + 64ROT + 128SBOX
}
\]

\subsubsection{加密长度为 \(n\) 字节数据的复杂度}

设明文长度为 \(n\) 字节，由于 SM4 每个分组大小为 16 字节，所以需要处理的分组数为

\[
b=\left\lceil \frac{n}{16} \right\rceil
\]
如果密钥只扩展一次，那么总操作量可以表示为

\[
T(n)
\approx
T_{\text{key}} + b \cdot T_{\text{block}}
\]
即
\[
T(n)
\approx
(192XOR + 64ROT + 128SBOX)
+
\left\lceil \frac{n}{16} \right\rceil
(256XOR + 128ROT + 128SBOX)
\]
因此，从渐进复杂度角度看，SM4 对长度为 \(n\) 字节的数据进行加密时，总时间复杂度为
\[
\boxed{O(n)}
\]
但是更具体地说，SM4 每处理 16 字节数据，需要进行 32 轮迭代，总共大约包含
\[
\boxed{
256\text{ 次异或}
+
128\text{ 次循环移位}
+
128\text{ 次S盒查表}
}
\]
